% Generated draft: requires mathematical and source review.
\documentclass[11pt]{article}
\usepackage[margin=25mm]{geometry}
\usepackage{fontspec}
\setmainfont{Latin Modern Roman}
\usepackage{amsmath,amssymb,mathtools,unicode-math}
\setmathfont{Latin Modern Math}
\usepackage{enumitem,needspace,xurl,hyperref}
\setlength{\parindent}{0pt}
\setlength{\parskip}{5pt}
\setlength{\emergencystretch}{3em}
\newcommand{\N}{\mathbb{N}}
\newcommand{\Z}{\mathbb{Z}}
\newcommand{\Q}{\mathbb{Q}}
\newcommand{\R}{\mathbb{R}}
\newcommand{\C}{\mathbb{C}}
\newcommand{\F}{\mathbb{F}}
\newcommand{\A}{\mathbb{A}}
\newcommand{\PP}{\mathbb{P}}
\newcommand{\E}{\mathbb{E}}
\newcommand{\eps}{\varepsilon}
\newcommand{\dd}{\,\mathrm{d}}
\begin{document}
\section*{UnsolvedMath ID: 30002108}
\textbf{Source problem code:} catalogue entry
\subsection{Definitions and mathematical statement}
A parahoric level refers to a level condition at a prime \(p\) specified by a parahoric subgroup of the relevant reductive group over \(\Q_p\). The supplied record does not specify the prime, Shimura datum, remaining level structure, or hypotheses on the reductive group.

\textbf{Source question.} For every Shimura variety with parahoric level, does there exist an associated local-model diagram
\[
S_K \xleftarrow{\ \pi\ } \widetilde S_K
\xrightarrow{\ q\ } M^{\mathrm{loc}}
\]
such that \(\pi\) is a \(K\)-torsor and \(q\) is \(K\)-equivariant and smooth?

Here \(S_K\), \(M^{\mathrm{loc}}\), and the acting object denoted by \(K\) retain the source's notation. The record does not specify whether \(S_K\) denotes an integral model, the base over which the diagram is defined, the construction of the associated local model, or the interpretation of \(K\) as an object acting on schemes. In particular, the associated local model must be specified by the intended theory; allowing an arbitrary target scheme would not express the intended question.

To describe the torsor and equivariance requirements precisely once the missing data are supplied, suppose that a base scheme \(B\), a group scheme \(\mathcal G/B\), a \(B\)-scheme \(S_K\), and the intended \(\mathcal G\)-scheme \(M^{\mathrm{loc}}\) have been fixed. Under a left-action, fppf-torsor interpretation, the requirements would be that \(\pi\) is faithfully flat and locally of finite presentation and that
\[
\mathcal G\times_B\widetilde S_K
\longrightarrow
\widetilde S_K\times_{S_K}\widetilde S_K,
\qquad (g,t)\longmapsto(g\cdot t,t),
\]
is an isomorphism, while \(q\) is smooth and satisfies
\[
q(g\cdot t)=g\cdot q(t)
\]
functorially on points over every \(B\)-scheme. This explains one possible precise meaning of the displayed conditions; the supplied record does not establish that this group scheme, topology, or base is the intended one. A parahoric subgroup of a \(p\)-adic group and a group scheme used in a torsor condition are distinct kinds of objects. Thus the source question is preserved, but a fully specified conjecture requires these unresolved choices.
\subsection{Short English statement}
The question asks whether every Shimura variety with parahoric level has an associated local model reached by a smooth equivariant map from a torsor over its intended model. The source leaves the base, integral-model interpretation, acting group scheme, torsor topology, and local model unspecified, so the formulation requires human mathematical review.
\subsection{Sources}
\begin{itemize}
\item Supplied local UnsolvedMath registry record, permanent ID \(30002108\), source code \texttt{catalogue entry}, entitled \emph{Local-Model Diagrams for Parahoric Shimura Varieties}. The record gives proposed year 2012 and names no individual proposer.
\item Supplied report citation: \emph{Reductions of Shimura Varieties (2013)}, Oberwolfach Reports, DOI 10.4171/owr/2012/32, pp.\ 1961--2011. \url{https://doi.org/10.4171/owr/2012/32}. This is also the literature-sources entry labelled ``original OWR report.'' The report title/year and DOI are reproduced as supplied.
\item Supplied literature source: Kisin and Pappas, integral models with parahoric level. \url{https://pmihes.centre-mersenne.org/articles/10.1007/s10240-018-0100-0/}.
\item Supplied literature source: Hamacher and Kim, mod \(p\) points at parahoric level. \url{https://arxiv.org/abs/2010.10496}.
\item Literature assessment dated exactly \texttt{2026-08-21}: the supplied curation labels the problem ``partially solved'' and reports constructions of integral models and local-model diagrams for broad Hodge- and abelian-type classes with parahoric level, Kottwitz nearby-cycles statements under additional hypotheses, and a fully general conjectural diagram that remains conditional in the literature. This is supplied dated curation, not an independently verified assessment. The background explicitly requests independent verification of time-sensitive status and cited claims.
\item The registry labels the statement \texttt{corrected\_verified} and reports that its mathematical content was checked against the official Oberwolfach report, with notation, context, and question form normalized. This is a supplied verification claim; no source consultation or independent verification is asserted here.
\end{itemize}
\end{document}
